\begin{pastexam}[label={px:ncku-107-cce-06}]{成大電通所 107 學年度 問答題第 2 題}
（問答；10 分）

\textbf{2. Matrix} (10 pts)
\begin{enumerate}[label=\alph*.]
\item (\textbf{5 pts}) A matrix is said to be \textbf{lower triangular} when all elements below the main diagonal of square matrix is zero. Given a lower triangle \textbf{A}, each non-zero element is stored by rows in a one-dimensional array \textbf{B} with $a$[0][0] being stored in $b$[0] (see Fig. (a)). Give a single expression for $k$ in terms of $i$ and $j$ if an element $a_{i,j}$ is stored at $b[k]$.
\item (\textbf{5 pts}) A matrix is said to be a \textbf{symmetric matrix} if it is a square matrix and its transpose is equal to itself. Similarly, each element of a symmetric matrix is stored by rows in a one-dimensional array \textbf{B} with $a$[0][0] being stored in $b$[0]. EX: a$_{10}$ =a$_{01}$ is stored in b(1); a$_{11}$ is stored in b(2); a$_{21}$ =a$_{12}$ is stored in b(3) and so on (see Fig. (b)). Write a single expression for $k$ in terms of $i$ and $j$ if an element $a_{i,j}$ is stored at $b[k]$ (\textbf{hint:} the functions Max($i$,$j$) and Min($i$,$j$) could be used in the expression).
\end{enumerate}

\begin{center}
\begin{tikzpicture}[font=\small,>={Stealth[length=1.8mm]}]
\begin{scope}[xshift=0.00cm]
\draw (0.00,0)--(0.00,-3.00); \draw (0,0.00)--(3.00,0.00);
\draw (0.60,0)--(0.60,-3.00); \draw (0,-0.60)--(3.00,-0.60);
\draw (1.20,0)--(1.20,-3.00); \draw (0,-1.20)--(3.00,-1.20);
\draw (1.80,0)--(1.80,-3.00); \draw (0,-1.80)--(3.00,-1.80);
\draw (2.40,0)--(2.40,-3.00); \draw (0,-2.40)--(3.00,-2.40);
\draw (3.00,0)--(3.00,-3.00); \draw (0,-3.00)--(3.00,-3.00);
\node at (-0.4,-0.30) {0};
\node at (0.30,-0.30) {b$_{0}$};
\node at (0.90,-0.30) {0};
\node at (1.50,-0.30) {0};
\node at (2.10,-0.30) {0};
\node at (-0.4,-0.90) {1};
\node at (0.30,-0.90) {b$_{1}$};
\node at (0.90,-0.90) {b$_{2}$};
\node at (1.50,-0.90) {0};
\node at (2.10,-0.90) {0};
\node at (-0.4,-1.50) {2};
\node at (0.30,-1.50) {b$_{3}$};
\node at (0.90,-1.50) {b$_{4}$};
\node at (1.50,-1.50) {b$_{5}$};
\node at (2.10,-1.50) {0};
\node at (-0.4,-2.10) {3};
\node at (0.30,-2.10) {b$_{6}$};
\node at (0.90,-2.10) {b$_{7}$};
\node at (1.50,-2.10) {b$_{8}$};
\node at (2.10,-2.10) {b$_{9}$};
\node at (0.30,0.35) {0};
\node at (0.90,0.35) {1};
\node at (1.50,0.35) {2};
\node at (2.10,0.35) {3};
\draw[<-] (2.55,0.35)--++(0.45,0) node[right] {$j$};
\draw[<-] (-0.4,-2.55)--++(0,-0.4) node[below] {$i$};
\node at (1.50,-3.30) {(a)};
\end{scope}
\begin{scope}[xshift=5.00cm]
\draw (0.00,0)--(0.00,-3.00); \draw (0,0.00)--(3.00,0.00);
\draw (0.60,0)--(0.60,-3.00); \draw (0,-0.60)--(3.00,-0.60);
\draw (1.20,0)--(1.20,-3.00); \draw (0,-1.20)--(3.00,-1.20);
\draw (1.80,0)--(1.80,-3.00); \draw (0,-1.80)--(3.00,-1.80);
\draw (2.40,0)--(2.40,-3.00); \draw (0,-2.40)--(3.00,-2.40);
\draw (3.00,0)--(3.00,-3.00); \draw (0,-3.00)--(3.00,-3.00);
\node at (-0.4,-0.30) {0};
\node at (0.30,-0.30) {b$_{0}$};
\node at (0.90,-0.30) {b$_{1}$};
\node at (1.50,-0.30) {b$_{3}$};
\node at (2.10,-0.30) {b$_{6}$};
\node at (-0.4,-0.90) {1};
\node at (0.30,-0.90) {b$_{1}$};
\node at (0.90,-0.90) {b$_{2}$};
\node at (1.50,-0.90) {b$_{4}$};
\node at (2.10,-0.90) {b$_{7}$};
\node at (-0.4,-1.50) {2};
\node at (0.30,-1.50) {b$_{3}$};
\node at (0.90,-1.50) {b$_{4}$};
\node at (1.50,-1.50) {b$_{5}$};
\node at (2.10,-1.50) {b$_{8}$};
\node at (-0.4,-2.10) {3};
\node at (0.30,-2.10) {b$_{6}$};
\node at (0.90,-2.10) {b$_{7}$};
\node at (1.50,-2.10) {b$_{8}$};
\node at (2.10,-2.10) {b$_{9}$};
\node at (0.30,0.35) {0};
\node at (0.90,0.35) {1};
\node at (1.50,0.35) {2};
\node at (2.10,0.35) {3};
\draw[<-] (2.55,0.35)--++(0.45,0) node[right] {$j$};
\draw[<-] (-0.4,-2.55)--++(0,-0.4) node[below] {$i$};
\node at (1.50,-3.30) {(b)};
\end{scope}
\end{tikzpicture}
\end{center}\end{pastexam}

\begin{pastsoln}
這份考卷和成大電機所 107 學年度的「資料結構」是同一份（兩系所共用），完整詳解見歷屆試題~\ref{px:ncku-107-ee-ds-06}。

答案：a. $k = \dfrac{i(i+1)}{2} + j$（$i \ge j$）；b. $k = \dfrac{\mathrm{Max}(i,j)\,(\mathrm{Max}(i,j)+1)}{2} + \mathrm{Min}(i,j)$。
原題 a 小題的定義把「以上」寫成「以下」，照題名和圖 (a) 作答。
\end{pastsoln}
